
\documentclass[../../ms_q.tex]{subfiles}

\begin{document}

\setlength{\abovedisplayskip}{2mm}
\setlength{\belowdisplayskip}{1mm}

\section{数学Ⅱ～式と証明～}

\subsection{多項式の割り算}

\begin{enumerate}
    \item 整式 $A = 2x^3 - 5x^2 + x + 2$ を整式 $B$ で割ると，
        商が $2x - 1$，余りが $x + 1$ であった．整式 $B$ を求めよ．
    
    \item 整式 $P(x) = x^3 + ax^2 + bx - 6$ を $x - 1$ で割ると
        余りが $-8$ となり，$x + 2$ で割ると割り切れるという．定数 $a, b$ の値を求めよ．
    
    \item 整式 $P(x)$ を $x - 2$ で割ると余りが $7$，
        $x + 1$ で割ると余りが $-2$ である．
        $P(x)$ を $(x - 2)(x + 1)$ で割ったときの余りを求めよ．
    
    \item 整式 $P(x)$ を $(x - 1)^2$ で割ると余りが $2x + 3$ であり，
        $x + 2$ で割ると余りが $11$ である．
        $P(x)$ を $(x - 1)^2(x + 2)$ で割ったときの余りを求めよ．
\end{enumerate}



\end{document}
